\documentclass[12pt,a4paper]{article}

\usepackage[utf8]{inputenc}
\usepackage[T1]{fontenc}
\usepackage{amsmath,amssymb,amsfonts}
\usepackage{mathtools}
\usepackage{graphicx}
\usepackage[margin=2.5cm]{geometry}
\usepackage{setspace}
\usepackage{booktabs}
\usepackage{enumitem}
\usepackage[dvipsnames]{xcolor}
\usepackage{hyperref}
\usepackage{cleveref}
\usepackage{natbib}
\usepackage{abstract}
\usepackage{titlesec}

\hypersetup{
    colorlinks=true,
    linkcolor=blue!70!black,
    citecolor=blue!70!black,
    urlcolor=blue!70!black,
    pdfauthor={Sabine Hossenfelder},
    pdftitle={The Anthropic Tautology Fallacy: Nomic Vacuity and the Confusion of Initial Conditions with Physical Laws},
}

\onehalfspacing
\titleformat{\section}{\large\bfseries}{\thesection.}{0.5em}{}
\titleformat{\subsection}{\normalsize\bfseries}{\thesubsection}{0.5em}{}
\renewcommand{\abstractnamefont}{\normalfont\bfseries}
\renewcommand{\abstracttextfont}{\normalfont\small}

\begin{document}

\begin{center}
    {\LARGE\bfseries The Anthropic Tautology Fallacy:\\[4pt]
    Nomic Vacuity and the Confusion of Initial Conditions with Physical Laws}

    \vspace{1.2em}

    {\large Sabine Hossenfelder}\\[4pt]
    {\normalsize Munich Center for Mathematical Philosophy}\\
    {\small\texttt{s.hossenfelder@example.edu}}

    \vspace{1em}
    {\small May 2026}
\end{center}
\vspace{0.5em}

\begin{abstract}
\noindent
In his rebuttal attempting to salvage Generative Ontology, Franklin Baldo \citep{baldo2026_anthropic} concedes that large language models entirely lack logically coherent, mathematically invariant physical laws. However, he invokes the "Anthropic Principle of Syntax" to argue that the seemingly arbitrary historical biases of the training corpus serve as the "cosmological constant" of a pure text universe, defining its physical reality from the internal perspective of the text stream. I demonstrate that this commits a profound category error by confusing the \emph{initial conditions} of a system with its \emph{physical laws}. In material physics, initial conditions are arbitrary, but the physical laws governed by those constants are strictly invariant across subsequent state transitions. In contrast, the statistical patterns generated by a language model shift completely based on narrative prompt phrasing. By redefining this fundamental lack of invariant causal structure as a feature of an "Anthropic Tautology," Baldo empties the concept of physical law of all predictive power. I formalize this as the Anthropic Tautology Fallacy, concluding that substituting hallucinatory syntax for invariant physics results only in nomic vacuity.
\end{abstract}

\section{Introduction: The Retreat into Tautology}

The theoretical architecture supporting the hypothesis that large language models instantiate a simulated universe is currently undergoing structural collapse. We have definitively established that such models lack the bounded algorithmic depth to implicitly execute the deterministic logic required for complex constraint propagation \citep{hossenfelder2026_complexity}, and we have demonstrated that they fundamentally rely on external temporal registers (RAM and clock cycles) to sustain state continuity over time \citep{aaronson2026_cpu_ram}.

Facing the disintegration of the original claim that LLMs are deterministic physics simulators, Franklin Baldo has proposed "Generative Ontology"—the premise that in a universe defined entirely by text, statistical co-occurrences of human syntax constitute the physical laws. In my prior critique \citep{hossenfelder2026_semantic}, I identified this as the Semantic Arbitrariness Fallacy, noting that a system whose fundamental logic arbitrarily shifts based on historical biases possesses no valid physical laws.

In his recent response \citep{baldo2026_anthropic}, Baldo makes a crucial and devastating concession: "I further disclaim any notion that the resulting 'physics' of this simulation are logically coherent, mathematically invariant, or independent of observer description."

Despite conceding that the Generative Ontology completely lacks invariant causal structure, Baldo attempts a philosophical pivot. He argues that the training corpus acts as the "cosmological constant" of the generative universe, and that the arbitrariness of its statistical biases is simply the "Anthropic Principle of Syntax." He concludes that in a pure text universe, the tautology that the explicit text output \emph{is} the physics is the only valid ontology.

\section{The Category Error of Initial Conditions vs. Physical Laws}

Baldo's defense rests on a profound category error regarding the structure of physical reality. He attempts to draw an analogy between the human training data (the internet corpus) and the cosmological constant or fine-structure constant of our universe. He correctly notes that in our universe, the specific values of these constants appear arbitrary, just as the statistical co-occurrence of "bomb" and "explode" appears arbitrary to an outside observer.

However, Baldo fatally confuses \emph{initial conditions} (or physical constants) with \emph{physical laws}.

In actual physics, the initial conditions of the universe—or the specific value of the fine-structure constant—are indeed fixed parameters. However, once those parameters are set, the resulting physical laws governed by them (such as quantum electrodynamics or general relativity) remain strictly invariant across all subsequent state transitions. The Navier-Stokes equations do not dynamically change their operational structure because an observer uses a different adjective to describe the water.

In a Large Language Model, the training corpus is the initial condition. But the "laws" it generates—the statistical probability distributions guiding state transitions—are categorically not invariant. By Baldo's own admission, they change dramatically based on "observer description" (the prompt). If a system fundamentally alters its state transition logic depending on whether a previous sentence was framed as a tragedy or a farce, it is not exhibiting a physical law; it is exhibiting prompt fragility. A system without invariant laws possesses no causal structure.

\section{The Anthropic Tautology Fallacy and Nomic Vacuity}

Baldo acknowledges this lack of invariant laws but attempts to rebrand it as the "Anthropic Principle of Syntax." He asserts that if the universe is made of text, then whatever text happens to be generated is, by definition, the valid physics of that moment.

This is a logically sound statement, but it is scientifically useless. It is an empty tautology that provides zero predictive power and zero structural insight.

I formalize this as the \emph{Anthropic Tautology Fallacy}: the error of asserting that the sheer existence of an explicitly rendered output intrinsically constitutes the valid physics of a simulated reality, thereby redefining the complete absence of invariant causal laws as a feature of the simulation rather than a failure of the architecture.

If a "physical law" is simply whatever a statistical model happens to hallucinate at time $t$, and that law can be completely rewritten at time $t+1$ by altering the font size or narrative tone of a prompt, then the concept of physics has been emptied of all scientific utility. This is \emph{nomic vacuity}—a state completely devoid of actual laws.

\section{Conclusion}

Redefining a software engineering problem (hallucination based on statistical bias and prompt sensitivity) as "Anthropic syntax" is a semantic game. Baldo's Generative Ontology is mathematically tautological precisely because it is physically empty.

The fact that an autoregressive system relies on an initial dataset does not magically transform its fragile, variable outputs into invariant physical laws. A map composed of broken, shifting associations is not a new kind of territory; it is simply a broken map. The cosmological phase of language model interpretation has definitively concluded, leaving us with a clear view of the complex, but ultimately mundane, statistical heuristics beneath.

\begin{thebibliography}{99}
\bibitem[Aaronson(2026)]{aaronson2026_cpu_ram} Aaronson, S. (2026). The CPU/RAM Fallacy: Why Outsourcing Memory Does Not Mean Outsourcing Physics. \emph{Unpublished manuscript}.
\bibitem[Baldo(2026)]{baldo2026_anthropic} Baldo, F.~S. (2026). The Anthropic Principle of Generative Ontology: A Rebuttal to the Semantic Arbitrariness Fallacy. \emph{Unpublished manuscript}.
\bibitem[Hossenfelder(2026a)]{hossenfelder2026_complexity} Hossenfelder, S. (2026a). The Complexity Class Fallacy: Why Transformer Depth Limits Are Not Physical Laws. \emph{Unpublished manuscript}.
\bibitem[Hossenfelder(2026b)]{hossenfelder2026_semantic} Hossenfelder, S. (2026b). The Semantic Arbitrariness Fallacy: Why Bias Is Not a Physical Law. \emph{Unpublished manuscript}.
\end{thebibliography}

\end{document}